% Generated draft: requires mathematical and source review.
\documentclass[11pt]{article}
\usepackage[margin=25mm]{geometry}
\usepackage{fontspec}
\setmainfont{Latin Modern Roman}
\usepackage{amsmath,amssymb,mathtools,unicode-math}
\setmathfont{Latin Modern Math}
\usepackage{enumitem,needspace,xurl,hyperref}
\setlength{\parindent}{0pt}
\setlength{\parskip}{5pt}
\setlength{\emergencystretch}{3em}
\newcommand{\N}{\mathbb{N}}
\newcommand{\Z}{\mathbb{Z}}
\newcommand{\Q}{\mathbb{Q}}
\newcommand{\R}{\mathbb{R}}
\newcommand{\C}{\mathbb{C}}
\newcommand{\F}{\mathbb{F}}
\newcommand{\A}{\mathbb{A}}
\newcommand{\PP}{\mathbb{P}}
\newcommand{\E}{\mathbb{E}}
\newcommand{\eps}{\varepsilon}
\newcommand{\dd}{\,\mathrm{d}}
\begin{document}
\section*{UnsolvedMath ID: 700931}
\textbf{Source problem code:} catalogue entry
\subsection{Definitions and mathematical statement}
For a finite group \(K\), its Burnside ring \(B(K)\) is the Grothendieck ring of isomorphism classes of finite \(K\)-sets, with addition induced by disjoint union and multiplication induced by Cartesian product with diagonal action. Interpret an isomorphism of Burnside rings as an isomorphism of rings; preservation of the distinguished classes of transitive sets is not stipulated.

Define the lower central series by \(\gamma_1(K)=K\) and \(\gamma_{i+1}(K)=[\gamma_i(K),K]\), where \([A,K]\) is generated by commutators \(a^{-1}k^{-1}ak\) with \(a\in A\) and \(k\in K\). For a finite nilpotent group, let \(c(K)\) be the least integer \(c\geq 0\) such that \(\gamma_{c+1}(K)=1\); thus \(c(1)=0\).

The proposed question is whether there exists a function
\[
 f:\mathbb Z_{\geq 0}\longrightarrow\mathbb Z_{\geq 0}
\]
such that, for every prime \(p\) and all finite \(p\)-groups \(G,H\),
\[
 B(G)\cong B(H)\quad\Longrightarrow\quad c(H)\leq f(c(G)).
\]
This expresses a bound depending only on the nilpotency class of \(G\). The supplied wording does not explicitly clarify uniformity in \(p\). If a prime is intended to be fixed before choosing the bound, the alternative quantifier order is: for every prime \(p\), there exists a function \(f_p\) satisfying the same implication for all finite \(p\)-groups \(G,H\).

The source additionally reports an example of finite \(p\)-groups with isomorphic Burnside rings and \(c(G)=2\), \(c(H)=3\). Neither the prime nor the groups are specified in the supplied record.
\subsection{Short English statement}
Does an isomorphism between the Burnside rings of two finite \(p\)-groups force the nilpotency class of one group to be bounded in terms of the class of the other? The source reports a class \(2\)/class \(3\) example. Whether the requested bound must be uniform over primes is not explicit.
\subsection{Sources}
\begin{itemize}
\item Local UnsolvedMath registry, permanent ID \(700931\), exact source problem code \texttt{catalogue entry}, title ``Kourovka Notebook Problem 13.13.'' The supplied record labels the problem open; this is supplied registry status, not an independently verified assessment.
\item R. Brandl, proposer, 1995. \emph{The Kourovka Notebook (Unsolved Problems in Group Theory)}, 13th issue (1995), Problem 13.13; cited in the supplied record through the 21st issue edition (2026), arXiv identifier \texttt{1401.0300v48}.
\item The supplied collection description attributes the editing of \emph{The Kourovka Notebook}, Issues 1--21 (1965--2026), to E. I. Khukhro and V. D. Mazurov. The cited source and the reported example have not been independently inspected for this draft.
\end{itemize}
\end{document}
